\chapter{Conditional Expectation, Rigorously}

In Chapter~\ref{ch37:expectation}, the Lebesgue integral gave us a single
formula for expectation, $\E[X] = \int_\Omega X\,\dd\P$, whether $X$ is
discrete, continuous, or mixed. Now we define conditional expectation
$\E[X \given \mathcal{G}]$ rigorously on sub-$\sigma$-algebras $\mathcal{G}$,
giving a definition that never divides by zero.

\section{Why we are doing this again}
\label{ch38:why-again}

In earlier chapters, conditional expectation was introduced for discrete variables
and continuous variables with joint densities: we wrote $g(x) = \E[Y \given X = x]$
and formed the random variable $\E[Y \given X] = g(X)$. We established properties
such as the tower property, the take-out rule, and the law of total variance.
Why revisit this topic with a measure-theoretic foundation?

Three key limitations in the elementary treatment motivate a more general formulation:

\begin{itemize}
  \item \textbf{Joint densities often do not exist.} The formula $\E[Y \given X = x] = \int y\,
    f_{Y \given X}(y \given x)\,\dd y$ requires that $(X,Y)$ possess a joint probability density.
    Many fundamental random objects --- such as Brownian motion paths, mixed discrete-continuous
    variables, and infinite-dimensional processes --- do not have joint densities with respect
    to Lebesgue measure.
  \item \textbf{Conditioning on null events requires careful definition.} The expression
    $\E[Y \given X = x]$ conditions on the event $\{X = x\}$, which has probability zero when $X$
    is continuous. As the Borel--Kolmogorov paradox demonstrates (Section~\ref{ch38:borel-kolmogorov}),
    conditioning on a set of measure zero without specifying a limiting family of conditioning
    events produces mathematical ambiguities.
  \item \textbf{Advanced processes condition on $\sigma$-algebras of information.}
    In martingale theory, one writes $\E[X_{n+1} \given \mathcal{F}_n] = X_n$, where
    $\mathcal{F}_n$ is the $\sigma$-algebra representing all information available up to time $n$.
    This requires conditioning directly on a $\sigma$-algebra rather than on the value of a
    single random variable.
\end{itemize}

Kolmogorov's resolution, from his 1933 \emph{Grundbegriffe}, is a striking change of tactics. He
stopped trying to \emph{construct} conditional expectation by a formula and instead
\emph{characterized} it by the one property it must have. The payoff is a single definition that
covers densities, masses, Brownian motion, and everything in between, and that turns out to be
exactly the geometric notion of \emph{orthogonal projection}. That equivalence is the heart of this
chapter, so we start there, with a picture you already trust.

\begin{focusbox}
This chapter assumes the measure-theoretic backbone from the previous two chapters: a probability
space $(\Omega, \mathcal{F}, \P)$, sigma-algebras as ``bundles of information,'' the Lebesgue
integral $\E[X] = \int_\Omega X \,\dd\P$, and the Radon--Nikodym theorem. If any of those feel hazy,
reread them first. Everything here is built on them.
\end{focusbox}

\section{The best-guess picture}
\label{ch38:best-guess}

Forget formulas for a moment. Conditional expectation answers a forecasting question: \emph{given
some information, what is my best guess for $Y$?} ``Best'' has to mean something, and the cleanest
meaning is mean-squared error. Among all guesses you are \emph{allowed} to make, which one minimizes
the expected squared miss?

What does ``allowed'' mean? If your information is the value of $X$, then a guess is allowed exactly
when it is a function of $X$ --- you cannot use information you do not have. So the question becomes:

\begin{equation}
\label{ch38:l2-problem}
  \text{minimize } \quad \E\!\big[(Y - g(X))^2\big] \quad \text{over all (measurable) functions } g.
\end{equation}

\begin{intuitionbox}
Imagine $Y$ is a person's height and $X$ is their shoe size. You will see many people with shoe size
$X = 9$; their heights scatter around some center. Your single allowed guess for all of them is
$g(9)$. To minimize squared error \emph{within that group}, you should aim at the group's mean
height. Do this for every shoe size and you have built the function $g(x) = \E[Y \given X = x]$. The
conditional expectation is just ``the within-group average, computed group by group.'' The whole
rigorous theory is a precise way to say exactly this when the groups are infinitesimally thin or do
not exist at all.
\end{intuitionbox}

Here is the calculus that confirms the intuition, before we abstract it. Fix a value $x$ and treat
$g(x)$ as a single number $c$ we get to choose. Conditionally on $X = x$,
\[
  \E\!\big[(Y - c)^2 \given X = x\big] = \Var(Y \given X = x) + \big(\E[Y \given X = x] - c\big)^2,
\]
by the same ``mean minimizes squared deviation'' identity you proved for plain expectation. The first
term does not depend on $c$; the second is a square, minimized (at zero) by $c = \E[Y \given X = x]$.
Minimizing pointwise in $x$ minimizes the average over $x$, so the best $g$ is exactly $g(x) =
\E[Y \given X = x]$. \emph{Conditional expectation is the least-squares predictor.}
This least-squares property provides the geometric foundation for the general definition.

\subsection{Geometric interpretation: orthogonal projection}
\label{ch38:projection-picture}

\begin{figure}[H]
\centering
\begin{tikzpicture}[scale=1.05, >=Stealth]
  % the subspace plane
  \fill[sourceblue!8] (-0.8,-0.3) -- (6.0,-0.3) -- (7.5,1.5) -- (0.7,1.5) -- cycle;
  \draw[sourceblue!60!black, thick] (-0.8,-0.3) -- (6.0,-0.3) -- (7.5,1.5) -- (0.7,1.5) -- cycle;
  \node[sourceblue!70!black, anchor=west, font=\small] at (6.8,0.2) {$L^2(\mathcal{G})$: allowed guesses};
  % origin
  \coordinate (O) at (0.6,0.2);
  \fill (O) circle (1.5pt);
  \node[below left, font=\small] at (O) {$0$};
  % Y point above plane
  \coordinate (Y) at (4.2,3.1);
  \fill[densityred] (Y) circle (1.8pt);
  \node[densityred, above=2pt, font=\small] at (Y) {$Y$};
  % projection foot
  \coordinate (P) at (4.2,0.55);
  \fill[targetgreen] (P) circle (1.8pt);
  \node[targetgreen, below=3pt, font=\small] at (P) {$\E[Y \given \mathcal{G}]$};
  % vectors
  \draw[->, thick, densityred, shorten >=2pt] (O) -- (Y) node[midway, above left, font=\small, densityred] {$Y$};
  \draw[->, thick, targetgreen, shorten >=2pt] (O) -- (P);
  \draw[->, thick, gray!70!black, shorten >=2pt] (P) -- (Y);
  \node[gray!70!black, font=\footnotesize, anchor=east] at (4.1,2.0) {$Y - \E[Y\given\mathcal{G}]$};
  % right angle mark
  \draw[gray!60!black] (P) ++(-0.22,0) -- ++(0,0.22) -- ++(0.22,0);
  % another candidate in plane
  \coordinate (Q) at (6.6,0.85);
  \fill (Q) circle (1.4pt);
  \node[right=2pt, font=\small] at (Q) {any other guess $g$};
  \draw[dashed, gray] (Q) -- (Y);
\end{tikzpicture}
\caption{Conditional expectation as orthogonal projection. The plane is the space of guesses you are
allowed to make (functions of your information $\mathcal{G}$). The random variable $Y$ lives above it.
The closest point in the plane --- the foot of the perpendicular --- is $\E[Y \given \mathcal{G}]$.
The error vector $Y - \E[Y \given \mathcal{G}]$ is orthogonal to the entire plane. Any other guess
$g$ is farther from $Y$ (Pythagoras).}
\label{ch38:fig-projection}
\end{figure}

If Figure \ref{ch38:fig-projection} looks like the picture of projecting a vector onto a
subspace from linear algebra, that is because it \emph{is} that picture. Making the
analogy exact is the cleanest way to understand existence and uniqueness at once.

\begin{connectionbox}
In the companion linear algebra book you projected a vector $\mathbf{v} \in \R^n$ onto a subspace
$W$. The projection $\hat{\mathbf{v}}$ is the unique point of $W$ closest to $\mathbf{v}$, and it is
characterized by orthogonality: $\mathbf{v} - \hat{\mathbf{v}} \perp \mathbf{w}$ for every
$\mathbf{w} \in W$, where $\perp$ means the inner product $\langle \cdot, \cdot \rangle$ vanishes.

The probabilistic dictionary:
\begin{center}
\begin{tabular}{ll}
\toprule
Linear algebra & Probability \\
\midrule
vectors $\mathbf{v}$ & random variables $Y$ with $\E[Y^2] < \infty$ \\
inner product $\langle \mathbf{u},\mathbf{v}\rangle$ & $\langle U, V \rangle = \E[UV]$ \\
length$^2$ $\norm{\mathbf{v}}^2$ & $\E[Y^2]$ \\
subspace $W$ & $\mathcal{G}$-measurable variables $L^2(\mathcal{G})$ \\
projection onto $W$ & conditional expectation $\E[\,\cdot \given \mathcal{G}]$ \\
$\mathbf{v}-\hat{\mathbf{v}} \perp W$ & error is uncorrelated with every $\mathcal{G}$-variable \\
\bottomrule
\end{tabular}
\end{center}
The space $L^2$ of square-integrable random variables, with inner product
$\langle U, V\rangle = \E[UV]$, is a complete inner-product space --- a \emph{Hilbert space}.
The Hilbert-space projection theorem guarantees that onto any closed subspace there exists
a unique nearest point. Conditional expectation is precisely this orthogonal projection;
existence and uniqueness in $L^2$ follow directly from the Hilbert-space projection theorem.
\end{connectionbox}

\section{Conditioning on a sigma-algebra: the definition}
\label{ch38:definition}

We now make the leap that the elementary theory could not. Instead of conditioning on a value or even
on a variable, we condition on a sigma-algebra $\mathcal{G} \subseteq \mathcal{F}$. Think of
$\mathcal{G}$ as the collection of all events whose outcome you are allowed to know --- your
information.

\begin{intuitionbox}
A sigma-algebra is a precise way to say ``how much can you resolve.'' If $\mathcal{G}$ is generated
by a variable $X$, then knowing $\mathcal{G}$ means knowing $X$: you can tell apart two outcomes
exactly when they give different $X$ values. A coarse $\mathcal{G}$ blurs many outcomes together; the
finest $\mathcal{G} = \mathcal{F}$ resolves everything. ``Conditioning on $\mathcal{G}$'' means
``give the best forecast computable from the information $\mathcal{G}$.''
\end{intuitionbox}

Two demands pin down conditional expectation. First, your forecast must be \emph{computable from your
information}: it must be a $\mathcal{G}$-measurable random variable. Second, it must be \emph{right on
average over every event you can see}: over any $G \in \mathcal{G}$, the forecast and the true $Y$
must have the same integral. The second demand is the rigorous form of ``correct group averages.''

\begin{definition}[Conditional expectation given a sigma-algebra]
\label{ch38:def-cond-exp}
Let $(\Omega, \mathcal{F}, \P)$ be a probability space, $Y$ an integrable random variable
($\E\abs{Y} < \infty$), and $\mathcal{G} \subseteq \mathcal{F}$ a sub-sigma-algebra. A
\textbf{conditional expectation of $Y$ given $\mathcal{G}$} is any random variable, written
$\E[Y \given \mathcal{G}]$, satisfying:
\begin{enumerate}
  \item \textbf{(Measurability)} $\E[Y \given \mathcal{G}]$ is $\mathcal{G}$-measurable.
  \item \textbf{(Averaging / partial-integral identity)} For every $G \in \mathcal{G}$,
  \begin{equation}
  \label{ch38:defining-property}
    \int_G \E[Y \given \mathcal{G}] \,\dd\P \;=\; \int_G Y \,\dd\P,
    \qquad \text{equivalently} \qquad
    \E\!\big[\ind_G\,\E[Y\given\mathcal{G}]\big] = \E[\ind_G\, Y].
  \end{equation}
\end{enumerate}
\end{definition}

\begin{trapbox}
$\E[Y \given \mathcal{G}]$ is a \emph{random variable}, not a number. It is only defined ``almost
surely'': any two versions agree with probability one. The notation hides this, so be careful ---
statements about $\E[Y\given\mathcal{G}]$ are statements that hold a.s. Also note we never divide by
$\P(G)$ anywhere in \eqref{ch38:defining-property}. That division --- the source of all the trouble
with probability-zero conditioning --- has vanished. We compare \emph{integrals}, which behave
perfectly even when $\P(G) = 0$ pieces are involved.
\end{trapbox}

\subsection{Reading the definition as orthogonality}
\label{ch38:read-as-orthogonality}

The averaging property \eqref{ch38:defining-property} says $\E[\ind_G (Y - \E[Y\given\mathcal{G}])] =
0$ for all $G \in \mathcal{G}$. Indicators $\ind_G$ are the simplest $\mathcal{G}$-measurable
variables; by approximating any bounded $\mathcal{G}$-measurable $Z$ with simple functions and taking
limits, the property upgrades to
\begin{equation}
\label{ch38:orthogonality}
  \E\!\big[Z\,(Y - \E[Y \given \mathcal{G}])\big] = 0 \qquad \text{for every bounded }
  \mathcal{G}\text{-measurable } Z.
\end{equation}
In inner-product language: the error $Y - \E[Y\given\mathcal{G}]$ is orthogonal to every variable in
$L^2(\mathcal{G})$. That is precisely the orthogonality condition of Figure
\ref{ch38:fig-projection}. \textbf{Kolmogorov's definition and the projection picture are literally
the same equation.}

\section{Why it exists: two routes to the same object}
\label{ch38:existence}

A definition by axioms is only worth something if such an object actually exists, and is essentially
unique. We give the intuition for both. Uniqueness is short; existence has two proofs depending on
how much you assume about $Y$.

\subsection{Uniqueness (almost surely)}
\label{ch38:uniqueness}

\begin{proposition}[Uniqueness]
\label{ch38:prop-uniqueness}
If $V$ and $W$ both satisfy Definition \ref{ch38:def-cond-exp}, then $V = W$ almost surely.
\end{proposition}

\begin{whybox}
Both are $\mathcal{G}$-measurable and both integrate to $\int_G Y$ over every $G \in \mathcal{G}$, so
$\int_G (V - W)\,\dd\P = 0$ for all $G \in \mathcal{G}$. Now $D := V - W$ is $\mathcal{G}$-measurable.
Take $G = \set{D > \eps} \in \mathcal{G}$; then $0 = \int_G D \ge \eps\,\P(G)$, forcing
$\P(D > \eps) = 0$. Let $\eps \downarrow 0$ to get $\P(D > 0) = 0$; by symmetry $\P(D < 0) = 0$. So
$D = 0$ a.s. Testing the difference against the event $\set{D > \eps} \in \mathcal{G}$
is a standard technique in measure theory: it establishes that two conditional expectations
are almost surely equal without requiring explicit calculation of their values.
\end{whybox}

\begin{recoverybox}
\textbf{Pause and think:} Why did we test $D = V - W$ against the event $G = \{D > \eps\}$ rather than $\{D > 0\}$? Because on $\{D > \eps\}$, we have the strict lower bound $\int_G D\,\dd\P \ge \eps\,\P(D > \eps)$. Since $D$ is $\mathcal{G}$-measurable, $G \in \mathcal{G}$ is a valid test set, so the integral is 0, which immediately forces $\P(D > \eps) = 0$.
\end{recoverybox}

\subsection{Existence the easy way: the \texorpdfstring{$L^2$}{L2} projection}
\label{ch38:existence-l2}

Suppose first that $\E[Y^2] < \infty$. Then $Y$ is a vector in the Hilbert space $L^2(\Omega,
\mathcal{F}, \P)$, and $L^2(\mathcal{G})$ --- the square-integrable $\mathcal{G}$-measurable
variables --- is a \emph{closed} subspace of it. The Hilbert-space projection theorem hands us a
unique nearest point $\hat Y \in L^2(\mathcal{G})$, characterized by $\langle Y - \hat Y, Z\rangle =
0$ for all $Z \in L^2(\mathcal{G})$. Taking $Z = \ind_G$ recovers exactly
\eqref{ch38:defining-property}. So for square-integrable $Y$, conditional expectation simply
\emph{is} the orthogonal projection onto $L^2(\mathcal{G})$, and existence is the projection theorem.
This is the route worth remembering; it is clean, geometric, and complete.

\begin{happeningbox}
The reason the subspace must be \emph{closed} for the projection theorem to apply is the same reason
$L^2$ must be \emph{complete} (a Hilbert space, beyond a mere inner-product space): one takes limits.
Completeness of $L^2$ is itself a consequence of the dominated convergence theorem from the Lebesgue
integration chapter. The existence of conditional expectation in $L^2$ rests on this completeness.
\end{happeningbox}

\subsection{Existence the general way: Radon--Nikodym}
\label{ch38:existence-rn}

What if $Y$ is integrable but not square-integrable (so it is not a vector in $L^2$)? The projection
argument no longer applies directly. Kolmogorov's original construction handles every integrable $Y$
at once, using the Radon--Nikodym theorem. The idea is elegant: build a \emph{measure} out of $Y$,
notice it has a density with respect to $\P$, and \emph{that density is the conditional expectation.}

\begin{theorem}[Existence via Radon--Nikodym]
\label{ch38:thm-existence}
For any integrable $Y$ and any sub-sigma-algebra $\mathcal{G}$, a conditional expectation
$\E[Y\given\mathcal{G}]$ exists.
\end{theorem}

\begin{whybox}
Assume first $Y \ge 0$. Define a set function on the smaller sigma-algebra $\mathcal{G}$:
\[
  \nu(G) := \int_G Y \,\dd\P, \qquad G \in \mathcal{G}.
\]
This $\nu$ is a (finite) measure on $(\Omega, \mathcal{G})$: it is countably additive because the
integral is, by monotone convergence. The decisive link is that $\nu$ is \emph{absolutely continuous} with respect
to $\P$ restricted to $\mathcal{G}$: if $\P(G) = 0$ then $\int_G Y\,\dd\P = 0$, so $\nu(G) = 0$.
Whenever one measure is absolutely continuous with respect to another, the Radon--Nikodym theorem
produces a density --- a $\mathcal{G}$-measurable function $h \ge 0$ with
\[
  \nu(G) = \int_G h \,\dd\P \quad\text{for all } G \in \mathcal{G}.
\]
Compare with the definition of $\nu$: $\int_G h\,\dd\P = \int_G Y\,\dd\P$ for every $G \in
\mathcal{G}$, and $h$ is $\mathcal{G}$-measurable. That is precisely Definition
\ref{ch38:def-cond-exp}. So $h = \E[Y\given\mathcal{G}]$. For general integrable $Y$, split $Y = Y^+
- Y^-$ into positive and negative parts, apply the above to each, and subtract.
\end{whybox}

\begin{intuitionbox}
\textbf{The Radon--Nikodym derivative as local concentration.} What is the Radon--Nikodym derivative $\frac{\dd\nu}{\dd\P}$ conceptually? Think of $\P$ as the base volume of space (in cubic meters) and $\nu$ as the total mass of a substance distributed across that space (in kilograms). Absolute continuity $\nu \ll \P$ is the physical guarantee that vacuum carries no mass: if a region has zero volume ($\P(G) = 0$), it must contain zero mass ($\nu(G) = 0$). No point masses or infinitely dense singularities can exist in empty space. The Radon--Nikodym theorem states that under this condition, there is a local mass density function $h(\omega) = \frac{\dd\nu}{\dd\P}(\omega)$ (in $\text{kg}/\text{m}^3$) such that the total mass in any region $G$ is recovered by integrating density over volume: $\nu(G) = \int_G h(\omega)\,\dd\P(\omega)$. It is the exact multi-dimensional analogue of the Fundamental Theorem of Calculus: if $F(x) = \int_a^x f(t)\,\dd t$, then the density is the derivative $f(x) = F'(x) = \frac{\dd F}{\dd x}$.

When we restrict both measures to the coarser $\sigma$-algebra $\mathcal{G}$, the sensor can only measure mass and volume over the coarse pixel blocks of $\mathcal{G}$. The resulting Radon--Nikodym derivative $\frac{\dd\nu|_\mathcal{G}}{\dd\P|_\mathcal{G}}$ cannot see inside individual pixel blocks, so it must be constant across each block, equal to the average concentration $\frac{\nu(G)}{\P(G)} = \frac{\int_G Y\,\dd\P}{\P(G)}$ over that block. That coarse average density is precisely the conditional expectation $\E[Y\given\mathcal{G}]$.
\end{intuitionbox}

\begin{aestheticbox}
The Radon--Nikodym derivative $\dd\nu/\dd\P$ provides the general definition of
conditional expectation: $\E[Y\given\mathcal{G}]$ is a change-of-measure density.
This same analytical tool defines likelihood ratios in statistics and changes of
measure in mathematical finance (Girsanov's theorem), demonstrating how
measure-theoretic integration unifies diverse branches of probability.
\end{aestheticbox}

\begin{historybox}
Before Kolmogorov, conditioning on a continuous variable was computed by
dividing densities. The contradictions (such as the Borel--Kolmogorov paradox)
were well known: conditioning on the equator of a sphere gave different
conditional densities depending on the coordinate slices used to approach it.
In his 1933 monograph \emph{Grundbegriffe der Wahrscheinlichkeitsrechnung},
Andrey Kolmogorov defined conditional expectation via its partial-integral
averaging property, using the Radon--Nikodym theorem (proved by Otton Nikodym in
1930, building on Johann Radon's 1913 work) for existence. This definition
gave a rigorous base for Joseph Doob's martingale theory, Kiyosi It\^o's
stochastic calculus, and the Kalman filter.
\end{historybox}

\section{It really does recover everything you knew}
\label{ch38:recovers}

A new definition is only trustworthy if the old correct answers fall out of it. They do.

\subsection{Conditioning on a discrete partition}
\label{ch38:discrete-recover}

Let $\mathcal{G} = \sigma(B_1, B_2, \dots)$ be generated by a countable partition of $\Omega$ into
events $B_i$ with $\P(B_i) > 0$. A $\mathcal{G}$-measurable variable must be constant on each $B_i$
(it cannot distinguish points inside a block), so $\E[Y\given\mathcal{G}] = \sum_i c_i \ind_{B_i}$
for some constants $c_i$. The averaging property over $G = B_i$ reads $c_i\,\P(B_i) = \int_{B_i}
Y\,\dd\P$, hence
\[
  c_i = \frac{1}{\P(B_i)} \int_{B_i} Y \,\dd\P = \E[Y \given B_i].
\]
So on the block $B_i$, the conditional expectation equals the elementary $\E[Y \given B_i]$ --- the
within-group average from the height-and-shoe-size story. The new definition has reproduced the old
one exactly, with no density assumption.

\subsection{Conditioning on a continuous variable}
\label{ch38:continuous-recover}

Let $\mathcal{G} = \sigma(X)$ and suppose $(X, Y)$ has a joint density. A theorem (the Doob--Dynkin
lemma) says every $\sigma(X)$-measurable variable is a function of $X$, so $\E[Y\given\sigma(X)] =
g(X)$ for some measurable $g$. Plug into the averaging property over $G = \set{X \in A}$:
\[
  \int_{\set{X\in A}} g(X)\,\dd\P = \int_{\set{X\in A}} Y\,\dd\P
  \;\Longrightarrow\;
  \int_A g(x) f_X(x)\,\dd x = \int_A \!\!\int y\, f_{X,Y}(x,y)\,\dd y\,\dd x.
\]
Since this holds for every Borel set $A$, the integrands match for almost every $x$:
\[
  g(x) f_X(x) = \int y\, f_{X,Y}(x,y)\,\dd y
  \;\Longrightarrow\;
  g(x) = \int y\,\frac{f_{X,Y}(x,y)}{f_X(x)}\,\dd y = \int y\, f_{Y\given X}(y\given x)\,\dd y.
\]
The dividing-densities formula re-emerges --- but now as a \emph{consequence}, valid where the
densities exist, rather than as a shaky starting point. We write $\E[Y\given\sigma(X)]$ and
$\E[Y\given X]$ interchangeably; they mean the same random variable.

\begin{walkbox}
Convince yourself why the constant ``$\E[Y]$'' is the right answer for the trivial sigma-algebra
$\mathcal{G} = \set{\emptyset, \Omega}$. A $\mathcal{G}$-measurable variable must be constant
(it cannot resolve any event), say $\equiv c$. The averaging property over $G = \Omega$ gives
$c = \int_\Omega Y\,\dd\P = \E[Y]$. With \emph{no} information, your best guess is the overall mean.
The projection picture agrees: $L^2(\mathcal{G})$ is the line of constants, and the nearest constant
to $Y$ is $\E[Y]$.
\end{walkbox}

\section{The properties, re-proved from the definition}
\label{ch38:properties}

Every property you used in the elementary chapter now follows cleanly from measurability plus the
averaging identity. We prove the most important ones; the pattern is always the same --- propose a
candidate, check it is $\mathcal{G}$-measurable, verify the partial-integral identity, and invoke
uniqueness.

\begin{proposition}[Core properties]
\label{ch38:prop-core}
Let $Y, Y_1, Y_2$ be integrable and $\mathcal{G} \subseteq \mathcal{F}$.
\begin{enumerate}
  \item \textbf{Linearity:} $\E[aY_1 + bY_2 \given \mathcal{G}] = a\,\E[Y_1\given\mathcal{G}] +
    b\,\E[Y_2\given\mathcal{G}]$.
  \item \textbf{Taking out what is known:} if $Z$ is $\mathcal{G}$-measurable and $ZY$ integrable,
    then $\E[ZY \given \mathcal{G}] = Z\,\E[Y\given\mathcal{G}]$.
  \item \textbf{Tower / smaller wins:} if $\mathcal{H} \subseteq \mathcal{G}$, then
    $\E\big[\E[Y\given\mathcal{G}] \,\big|\, \mathcal{H}\big] = \E[Y \given \mathcal{H}]$.
  \item \textbf{Independence drops:} if $Y$ is independent of $\mathcal{G}$, then
    $\E[Y\given\mathcal{G}] = \E[Y]$.
  \item \textbf{Monotonicity \& Jensen:} $Y_1 \le Y_2$ a.s.\ implies $\E[Y_1\given\mathcal{G}] \le
    \E[Y_2\given\mathcal{G}]$ a.s.; and $\varphi(\E[Y\given\mathcal{G}]) \le \E[\varphi(Y)\given
    \mathcal{G}]$ for convex $\varphi$.
\end{enumerate}
\end{proposition}

\begin{whybox}[Taking out what is known]
We want to show $Z\,\E[Y\given\mathcal{G}]$ is the conditional expectation of $ZY$. It is
$\mathcal{G}$-measurable (product of two $\mathcal{G}$-measurable things). For the averaging property,
we must check $\int_G Z\,\E[Y\given\mathcal{G}]\,\dd\P = \int_G ZY\,\dd\P$ for every $G \in
\mathcal{G}$. Start with $Z = \ind_B$, $B \in \mathcal{G}$: the left side is $\int_{G\cap B}
\E[Y\given\mathcal{G}]\,\dd\P$, which by the defining property of $\E[Y\given\mathcal{G}]$ (applied to
the set $G \cap B \in \mathcal{G}$) equals $\int_{G\cap B} Y\,\dd\P = \int_G \ind_B Y\,\dd\P$. So it
holds for indicators; by linearity it holds for simple $\mathcal{G}$-measurable $Z$; by monotone
convergence it holds for all $\mathcal{G}$-measurable $Z \ge 0$, then for general $Z$ by splitting.
Uniqueness finishes the proof. The intuition: variables you already know act like \emph{constants}
once you have conditioned --- there is nothing left to average about them.
\end{whybox}

\subsection{The tower property, re-proved}
\label{ch38:tower}

The tower property is the engine of martingale theory, so it deserves a clean proof from scratch.

\begin{theorem}[Tower property]
\label{ch38:thm-tower}
If $\mathcal{H} \subseteq \mathcal{G} \subseteq \mathcal{F}$ and $Y$ is integrable, then
\[
  \E\big[\,\E[Y\given\mathcal{G}] \,\big|\, \mathcal{H}\,\big] = \E[Y\given\mathcal{H}]
  \qquad \text{a.s.}
\]
\end{theorem}

\begin{whybox}
Let $W := \E[Y\given\mathcal{G}]$ and let us show $\E[W \given \mathcal{H}] = \E[Y\given\mathcal{H}]$.
By uniqueness it suffices to show that $\E[Y\given\mathcal{H}]$ \emph{itself} qualifies as the
conditional expectation of $W$ given $\mathcal{H}$. It is $\mathcal{H}$-measurable by construction.
For the averaging property, take any $H \in \mathcal{H}$. Because $\mathcal{H} \subseteq \mathcal{G}$,
we also have $H \in \mathcal{G}$, so we may use the defining property of $W = \E[Y\given\mathcal{G}]$
at the set $H$:
\[
  \int_H W \,\dd\P \;=\; \int_H Y \,\dd\P \;=\; \int_H \E[Y\given\mathcal{H}] \,\dd\P,
\]
where the last equality is the defining property of $\E[Y\given\mathcal{H}]$ at the same $H \in
\mathcal{H}$. Thus $\int_H W\,\dd\P = \int_H \E[Y\given\mathcal{H}]\,\dd\P$ for all $H \in
\mathcal{H}$, and $\E[Y\given\mathcal{H}]$ is $\mathcal{H}$-measurable, so it is the conditional
expectation of $W$. Done.
\end{whybox}

\begin{happeningbox}
The core of the proof is that $H \in \mathcal{H} \subseteq \mathcal{G}$, ensuring
that the averaging set is valid for both $\sigma$-algebras. The tower property
expresses the transitivity of averaging: taking a fine conditional expectation
and then averaging coarsely yields the same result as conditioning on the coarser
$\sigma$-algebra directly. In the projection framework (Figure~\ref{ch38:fig-projection}),
$L^2(\mathcal{H})$ is a closed subspace of $L^2(\mathcal{G})$, and projecting
first onto the larger subspace and then onto the smaller one is identical to
projecting directly onto the smaller subspace.
\end{happeningbox}

Setting $\mathcal{H} = \set{\emptyset, \Omega}$ recovers the law of total expectation
$\E\big[\E[Y\given\mathcal{G}]\big] = \E[Y]$, the special case you have been using all along.

\section{A picture: the conditional expectation as a step function}
\label{ch38:step-figure}

Let us see the projection concretely. Take $X \sim \Unif(0,1)$ and $Y = X^2 + \text{noise}$, so the
true regression curve is $\E[Y\given X = x] = x^2$. If instead of knowing $X$ exactly we only know
\emph{which quarter} $X$ falls in --- the coarse sigma-algebra $\mathcal{G}$ generated by the
partition $[0,\tfrac14), [\tfrac14,\tfrac12), [\tfrac12,\tfrac34), [\tfrac34,1)$ --- then
$\E[Y\given\mathcal{G}]$ must be constant on each quarter, equal to the average of $x^2$ over that
quarter. It is the best step-function approximation to $x^2$.

\begin{figure}[H]
\centering
\begin{tikzpicture}
\begin{axis}[
  width=11cm, height=6.6cm,
  title={Conditional expectation as orthogonal projection onto partition steps},
  title style={font=\small},
  xlabel={$x$}, ylabel={prediction of $Y$},
  domain=0:1, samples=120,
  xmin=0, xmax=1, ymin=0, ymax=1.05,
  legend pos=north west, legend cell align=left,
  grid=both, grid style={gray!18},
]
  % true regression curve E[Y|X=x] = x^2
  \addplot[densityred, thick, samples=120] {x^2};
  \addlegendentry{$\E[Y\given X]=x^2$ (full info)}
  % step function: average of x^2 over each quarter
  % avg over [a,b] of x^2 = (b^3-a^3)/(3(b-a))
  \addplot[targetgreen, thick, const plot mark left] coordinates {
    (0,0.020833) (0.25,0.145833) (0.5,0.395833) (0.75,0.770833) (1,0.770833)
  };
  \addlegendentry{$\E[Y\given\mathcal{G}]$ (quarter info)}
  % vertical partition lines
  \draw[gray!55, dashed] (axis cs:0.25,0) -- (axis cs:0.25,1.05);
  \draw[gray!55, dashed] (axis cs:0.5,0) -- (axis cs:0.5,1.05);
  \draw[gray!55, dashed] (axis cs:0.75,0) -- (axis cs:0.75,1.05);
\end{axis}
\end{tikzpicture}
\caption{Conditional expectation given coarser information. The smooth curve is $\E[Y\given X]=x^2$
(you know $X$ exactly). The step function is $\E[Y\given\mathcal{G}]$ when you know only which quarter
$X$ lies in: on each block it is the \emph{average} of the curve, the best constant fit there. As the
partition refines, the steps shrink toward the curve --- finer information, sharper forecast. This is
the projection of Figure \ref{ch38:fig-projection} drawn as a function.}
\label{ch38:fig-step}
\end{figure}

\begin{aestheticbox}
Figure~\ref{ch38:fig-step} illustrates the relationship between information,
subspaces, and approximation quality. Finer information corresponds to a richer
subspace of measurable functions, permitting a closer approximation to $Y$.
The tower property states that projecting a fine step function onto a coarser
subspace recovers the coarse step function directly. The unexplained variation ---
the expected squared difference between the predictor and the target --- corresponds
to the residual orthogonal component in $L^2$. This geometric perspective unifies
estimation, regression, and data-driven prediction within a single projection framework.
\end{aestheticbox}

\section{The paradox that forced this rigor}
\label{ch38:borel-kolmogorov}

Why insist on conditioning on sigma-algebras instead of just writing $\E[Y \given X = x]$? Because the
latter, taken literally, is \emph{ambiguous} when $\P(X = x) = 0$ --- and the ambiguity is not
academic.

\begin{paradoxbox}[Borel--Kolmogorov]
Pick a point uniformly at random on the surface of the Earth (a sphere). Now condition on the event
that the point lies \emph{on a particular great circle}. What is the conditional distribution of the
point along that circle?

\textbf{Answer 1.} Treat the great circle as a line of longitude (a meridian). The two-dimensional
uniform density on the sphere, expressed in latitude $\theta$ and longitude $\phi$, carries a factor
$\cos\theta$ (rings near the equator are bigger than rings near the poles). Conditioning on a fixed
longitude and normalizing the latitude, you get a distribution \emph{proportional to $\cos\theta$} ---
points bunch toward the equator.

\textbf{Answer 2.} Treat the same great circle as the equator (rotate your coordinates). Now by
symmetry the conditional distribution is \emph{uniform} along the circle.

Same circle, same uniform point on the sphere, two different conditional distributions. A
contradiction --- unless conditioning on a measure-zero event is not actually well defined by the
event alone.
\end{paradoxbox}

\begin{whybox}[Resolving it]
The resolution is exactly Kolmogorov's reframing. The event $\set{X = x}$ has probability zero, and
the symbol $\E[Y \given X = x]$ has \emph{no meaning in isolation}. What is well defined is
$\E[Y\given\sigma(X)]$, a single function whose values on probability-zero slices are pinned down only
up to a null set. To get a number on one particular slice you must say \emph{which family of thinning
events} converges to it --- meridians of longitude versus rings of latitude are two different limiting
families, and the Jacobian factor ($\cos\theta$ versus $1$) differs between them. The ``paradox'' is
the failure to specify the limiting procedure. Definition \ref{ch38:def-cond-exp} sidesteps it
entirely: it never conditions on a single null event, only ever integrates over positive-measure sets
$G \in \mathcal{G}$, where no ambiguity exists. The price is that the conditional expectation is
defined only a.s.; the reward is that it is defined \emph{consistently}.
\end{whybox}

\begin{trapbox}
The practical lesson: do not assign meaning to $\E[Y\given X = x]$ at an isolated point without a
density or a stated limiting scheme. With a joint density, the dividing-densities formula gives the
canonical answer (implicitly fixing the limiting family to Cartesian slices). Without
one, ask ``conditional with respect to which sigma-algebra?'' before you compute anything.
\end{trapbox}

\section{Why this is the right notion for martingales}
\label{ch38:martingales}

Now the payoff. A \textbf{filtration} is an increasing chain $\mathcal{F}_0 \subseteq \mathcal{F}_1
\subseteq \cdots$ of sigma-algebras: $\mathcal{F}_n$ is everything known by time $n$, and information
only accumulates. A process $(M_n)$ is a \textbf{martingale} (adapted to the filtration) if each $M_n$
is $\mathcal{F}_n$-measurable, integrable, and
\begin{equation}
\label{ch38:martingale-def}
  \E[M_{n+1} \given \mathcal{F}_n] = M_n \qquad \text{for all } n.
\end{equation}

This definition is \emph{impossible to state} with the elementary $\E[Y\given X = x]$ machinery: the
``information at time $n$'' is a whole sigma-algebra, generated by the entire past $M_0, \dots, M_n$
(and possibly more), not a single value. Conditioning on a sigma-algebra is not a luxury here; it is
the only language in which a martingale can be written down.

\begin{figure}[H]
\centering
\begin{tikzpicture}[scale=1.0, >=Stealth]
  % nested sigma-algebras as growing boxes (draw largest first so fills do not cover text)
  \draw[sourceblue!80!black, thick, rounded corners, fill=sourceblue!5] (0,0) rectangle (5.6,2.0);
  \draw[sourceblue!80!black, thick, rounded corners, fill=sourceblue!10] (0,0) rectangle (3.8,2.0);
  \draw[sourceblue!80!black, thick, rounded corners, fill=sourceblue!15] (0,0) rectangle (2.0,2.0);

  % labels on top of boxes
  \node[sourceblue!80!black, font=\small\bfseries] at (1.0,2.3) {$\mathcal{F}_0$};
  \node[align=center, font=\small] at (1.0,1.0) {known\\at $0$};

  \node[sourceblue!80!black, font=\small\bfseries] at (2.9,2.3) {$\mathcal{F}_1$};

  \node[sourceblue!80!black, font=\small\bfseries] at (4.7,2.3) {$\mathcal{F}_2$};
  \node[align=center, font=\small] at (4.7,1.0) {more\\known\\later};

  \node[gray!60!black, font=\large] at (6.5,1.0) {$\subseteq\cdots$};

  % the conditional expectation arrow
  \draw[->, very thick, densityred] (8.0,2.0) -- (8.0,0.0);
  \node[densityred, right=4pt, font=\small] at (8.0,2.0) {$M_{n+1}$ (random)};
  \fill[targetgreen] (8.0,0.0) circle (2.2pt);
  \node[targetgreen, right=4pt, font=\small] at (8.0,0.0) {$\E[M_{n+1}\given\mathcal{F}_n]=M_n$};
  \node[gray!60!black, right=4pt, align=left, font=\footnotesize] at (8.0,1.0)
    {project onto\\information\\known by time $n$};
\end{tikzpicture}
\caption{A filtration is a growing stack of information: $\mathcal{F}_0 \subseteq \mathcal{F}_1
\subseteq \cdots$. The martingale condition projects the unknown future value $M_{n+1}$ onto the
information $\mathcal{F}_n$ available now; for a fair game the projection lands back on the current
value $M_n$. Each step is one application of Definition \ref{ch38:def-cond-exp}.}
\label{ch38:fig-filtration}
\end{figure}

\begin{applicationbox}[The tower property is the fairness of a fair game]
Read \eqref{ch38:martingale-def} as: given all the history up to now, the expected next value equals
the current value --- a fair game, no drift (Figure \ref{ch38:fig-filtration}). The tower property
(Theorem \ref{ch38:thm-tower}) immediately gives the long-range version. For any $m < n$,
\[
  \E[M_n \given \mathcal{F}_m]
  = \E\big[\E[M_n\given\mathcal{F}_{n-1}] \,\big|\, \mathcal{F}_m\big]
  = \E[M_{n-1}\given\mathcal{F}_m] = \cdots = M_m,
\]
peeling off one step at a time via the tower. A fair game is fair over any horizon, far beyond a single step.
And taking $m=0$ with $\mathcal{F}_0$ trivial, $\E[M_n] = \E[M_0]$ for all $n$: the expected fortune
never changes. This single consequence underlies optional stopping, the gambler's-ruin computations,
and the no-arbitrage pricing of derivatives, where the discounted asset price is a martingale under
the risk-neutral measure.
\end{applicationbox}

\begin{happeningbox}
Why must the \emph{whole} sigma-algebra $\mathcal{F}_n$ appear, rather than just $M_n$? Because
fairness must hold given \emph{everything} a gambler could know, not merely the current price.
A process can satisfy $\E[M_{n+1}\given M_n] = M_n$ (fair given only the last value) yet fail
$\E[M_{n+1}\given\mathcal{F}_n] = M_n$ (using the full past lets you detect a pattern and beat the
game). The sigma-algebra formulation closes that loophole. ``You cannot beat a martingale'' is a
statement about \emph{all} usable information, which is exactly what a sigma-algebra encodes.
\end{happeningbox}

\section{Why this is the right notion for Bayesian updating}
\label{ch38:bayes}

Bayesian inference is conditioning made into a worldview: you hold a prior over an unknown $\Theta$,
observe data $D$, and update to the posterior. The posterior \emph{mean} $\E[\Theta \given D]$ is a
conditional expectation in exactly our sense, and the projection picture explains why it is the
optimal estimate.

\begin{applicationbox}[The posterior mean is the optimal point estimate]
Suppose you must report a single number $\hat\theta(D)$ to estimate $\Theta$, and you will be scored
by squared error $\E[(\Theta - \hat\theta(D))^2]$. Your estimate may depend only on the data, i.e.\ it
must be $\sigma(D)$-measurable. This is precisely the least-squares problem
\eqref{ch38:l2-problem} with $Y = \Theta$ and information $\mathcal{G} = \sigma(D)$. Its solution is
the projection:
\[
  \hat\theta(D) = \E[\Theta \given D] = \text{posterior mean}.
\]
\emph{Under squared-error loss, the Bayes-optimal estimator is the posterior mean} --- not by
convention, but because it is the orthogonal projection of the truth onto the space of data-driven
guesses. The data filtration grows as observations arrive; the tower property says yesterday's
posterior is the prior for today, and updating in one batch or sequentially gives the same answer.
Bayesian learning \emph{is} a sequence of projections onto ever-finer information.
\end{applicationbox}

\begin{connectionbox}
The same projection idea, restricted to \emph{linear} functions of the data instead of all functions,
gives linear regression and the Kalman filter. Linear regression projects $Y$ onto the subspace
spanned by $\set{1, X_1, \dots, X_p}$; conditional expectation projects onto the much larger subspace
of \emph{all} functions of the data. So $\E[Y\given X]$ is the ``true'' regression function and
ordinary least squares is its best linear approximation. When $(Y, X)$ is jointly Gaussian the two
coincide --- the conditional expectation is already linear --- which is the secret reason the Kalman
filter, built from linear projections, is exactly optimal in the Gaussian world. See the linear
algebra companion for the projection machinery and the estimation chapters here for the statistics.
\end{connectionbox}

\section{A fully worked example}
\label{ch38:worked}

\begin{example}[Conditioning a sum on one of its parts, three ways]
\label{ch38:ex-sum}
Let $X, Y$ be independent with $\E[X] = \mu_X$, $\E[Y] = \mu_Y$, and let $S = X + Y$. We compute
$\E[X \given S]$, $\E[S \given X]$, and verify the tower property, using only the rules of
Proposition \ref{ch38:prop-core}.

\textbf{(a) $\E[S\given X]$.} Write $S = X + Y$ and use linearity: $\E[S\given X] = \E[X\given X] +
\E[Y\given X]$. Now $X$ is $\sigma(X)$-measurable so $\E[X\given X] = X$ (take out what is known, with
$Y=1$). And $Y \indep X$, so $\E[Y\given X] = \E[Y] = \mu_Y$ (independence drops). Hence
\[
  \E[S\given X] = X + \mu_Y.
\]
Sensible: knowing $X$, the only randomness left in $S$ is $Y$, whose average is $\mu_Y$.

\textbf{(b) $\E[X\given S]$ when $X, Y$ are i.i.d.} Now suppose $X$ and $Y$ have the \emph{same}
distribution. By symmetry $\E[X\given S]$ and $\E[Y\given S]$ are a.s.\ equal as functions of $S$ (the
symmetric roles give identical $\sigma(S)$-measurable answers). Their sum is forced:
\[
  \E[X\given S] + \E[Y\given S] = \E[X + Y\given S] = \E[S\given S] = S,
\]
so each equals $S/2$:
\[
  \boxed{\;\E[X\given S] = \tfrac{S}{2}.\;}
\]
No densities, no integrals --- pure structure. If $X,Y \sim \Pois(\lambda)$ then given $S = n$, $X$ is
$\Bin(n, \tfrac12)$ with mean $n/2$, matching. If $X,Y \sim \Normal(0,\sigma^2)$, then given $S$, $X$
is Gaussian with mean $S/2$, matching again. One symmetry argument, every distribution at once.

\textbf{(c) Tower check.} We verify $\E\big[\E[X\given S]\big] = \E[X]$ in the i.i.d.\ case:
$\E[S/2] = \tfrac12(\mu_X + \mu_Y) = \mu_X$ since $\mu_X = \mu_Y$. The law of total expectation holds.
\end{example}

\begin{labbox}
Redo Example \ref{ch38:ex-sum}(b) without symmetry, when $X \sim \Pois(\alpha)$, $Y \sim
\Pois(\beta)$ are independent with $\alpha \ne \beta$. Show that given $S = n$, $X \sim \Bin(n,
\tfrac{\alpha}{\alpha+\beta})$ and therefore $\E[X\given S] = \tfrac{\alpha}{\alpha+\beta}\,S$. Notice
the answer is again \emph{linear} in $S$ even though no Gaussian is in sight --- and check it reduces
to $S/2$ when $\alpha = \beta$.
\end{labbox}

\section{Summary}
\label{ch38:summary}

\begin{recoverybox}
\begin{itemize}
  \item Conditional expectation $\E[Y\given\mathcal{G}]$ is the unique (a.s.)\ random variable that is
    (i) $\mathcal{G}$-measurable and (ii) has the same integral as $Y$ over every $G \in \mathcal{G}$.
    No division by probabilities, so it survives null events.
  \item Geometrically it is the \emph{orthogonal projection} of $Y$ onto $L^2(\mathcal{G})$ --- the
    best mean-squared predictor using the information $\mathcal{G}$. The error is uncorrelated with
    everything you know.
  \item Existence: for square-integrable $Y$, the Hilbert-space projection theorem; for any integrable
    $Y$, the Radon--Nikodym theorem, which makes $\E[Y\given\mathcal{G}]$ a change-of-measure density.
  \item It recovers the elementary definitions --- group averages over a partition, dividing densities
    for a continuous $X$, and $\E[Y]$ for no information --- as special cases.
  \item Linearity, taking out known factors, independence-drops, monotonicity, Jensen, and the
    \emph{tower property} all fall out by the candidate-and-uniqueness method. The tower property is
    transitivity of averaging, or projection-of-a-projection.
  \item The Borel--Kolmogorov paradox shows why conditioning on a bare null event is ill-posed; the
    sigma-algebra formulation is the cure.
  \item This is the native language of martingales ($\E[M_{n+1}\given\mathcal{F}_n] = M_n$) and of
    Bayesian updating (the posterior mean $\E[\Theta\given D]$ is the squared-error-optimal estimator).
\end{itemize}
\end{recoverybox}

\section{Epilogue: Looking Back and Looking Forward}
\label{ch38:epilogue}

With conditional expectation defined as an orthogonal projection onto a
sub-$\sigma$-algebra, we reach the end of the core chapters.

We started with counting and coin tosses, then built the Law of Large Numbers
and the Central Limit Theorem. From there we studied processes in time:
Markov chains, Poisson arrivals, random walks, and Brownian motion. In martingales,
we formalized fair games, and in Bayesian inference, we updated priors with data.

Whenever we pushed those elementary models to continuous spaces or infinite
sequences, formulas that divided by zero or swapped limits without justification
showed the need for clearer foundations. Measure theory resolved those issues:
random variables are measurable functions, expectations are Lebesgue integrals,
and conditional expectations are projections onto sub-$\sigma$-algebras.

The two appendices that follow gather reference material used throughout the text.
Appendix~\ref{ch39:sets} reviews sets, sums, and counting, while
Appendix~\ref{chAcla:calc-intro} collects multivariable calculus and linear algebra.

\section{Exercises}
\label{ch38:exercises}

\begin{exercise}[Quick mental checks]
\label{ch38:ex-warmup}
Warm up with three quick conceptual checks:
\begin{enumerate}
  \item If $\mathcal{G} = \{\emptyset, \Omega\}$ is the trivial $\sigma$-algebra containing zero information, what is $\E[Y \given \mathcal{G}]$? What is the orthogonal projection of $Y$ onto the subspace of constants?
  \item If $Y$ is already $\mathcal{G}$-measurable, what is $\E[Y \given \mathcal{G}]$? What is the error $Y - \E[Y \given \mathcal{G}]$?
  \item In Figure~\ref{ch38:fig-projection}, what is the geometric meaning of the identity $\E[(Y - \E[Y\given\mathcal{G}])\,Z] = 0$ for any bounded $\mathcal{G}$-measurable variable $Z$?
\end{enumerate}
\end{exercise}

\begin{exercise}[Definition drill]
\label{ch38:ex-def}
Let $\Omega = \set{1,2,3,4}$ with $\P$ uniform, and let $\mathcal{G} = \sigma(\set{1,2}, \set{3,4})$.
Let $Y(\omega) = \omega^2$. Compute $\E[Y\given\mathcal{G}]$ explicitly as a function on $\Omega$, and
verify the averaging property \eqref{ch38:defining-property} for $G = \set{1,2}$ and for $G =
\set{3,4}$.
\end{exercise}

\begin{exercise}[Smoothing reduces variance]
\label{ch38:ex-variance}
Assume $\E[Y^2] < \infty$. Using the orthogonality \eqref{ch38:orthogonality} with
$Z = \E[Y\given\mathcal{G}]$ (here a bounded approximation argument lets you take $Z$ square-integrable),
prove the
Pythagorean identity $\E[Y^2] = \E\big[\E[Y\given\mathcal{G}]^2\big] +
\E\big[(Y - \E[Y\given\mathcal{G}])^2\big]$. Deduce that $\Var(\E[Y\given\mathcal{G}]) \le \Var(Y)$:
conditioning never increases variance. When is it an equality?
\end{exercise}

\begin{exercise}[Take-out rule by hand]
\label{ch38:ex-takeout}
Let $X \sim \Unif(0,1)$ and $Y = X U$ where $U \sim \Unif(0,1)$ is independent of $X$. Find
$\E[Y\given X]$ using ``take out what is known'' plus ``independence drops,'' then verify your answer
by direct integration against $f_{Y\given X}$.
\end{exercise}

\begin{exercise}[Conditioning on a coarser sigma-algebra]
\label{ch38:ex-coarser}
For the setup of Figure \ref{ch38:fig-step} ($\E[Y\given X] = X^2$, $X \sim \Unif(0,1)$), compute the
four step heights of $\E[Y\given\mathcal{G}]$ for the quarter-partition by averaging $x^2$ over each
quarter, and confirm they match the values plotted. Then compute $\E\big[\E[Y\given\mathcal{G}]\big]$
two ways (from the steps, and as $\E[X^2] = 1/3$) to check the tower property.
\end{exercise}

\begin{exercise}[Tower in action]
\label{ch38:ex-tower}
Let $N \sim \Pois(\lambda)$ and, given $N = n$, let $S = \sum_{i=1}^n X_i$ with the $X_i$ i.i.d.\
mean $\mu$, independent of $N$. Using the tower property with $\mathcal{F} = \sigma(N)$, prove
$\E[S] = \lambda\mu$. Then identify the sigma-algebra you conditioned on and explain in one sentence
why $\E[S\given N] = N\mu$ qualifies as a conditional expectation per Definition
\ref{ch38:def-cond-exp}.
\end{exercise}

\begin{exercise}[Martingale from a fixed variable]
\label{ch38:ex-martingale}
Let $\mathcal{F}_0 \subseteq \mathcal{F}_1 \subseteq \cdots$ be a filtration and $Y$ an integrable
variable. Define $M_n := \E[Y\given\mathcal{F}_n]$ (a ``Doob martingale''). Use the tower property to
prove $(M_n)$ is a martingale, i.e.\ $\E[M_{n+1}\given\mathcal{F}_n] = M_n$. Interpret $M_n$ as ``the
current best estimate of $Y$ given what is known so far,'' and explain why a fair-game property is
exactly what you would expect of a well-behaved running estimate.
\end{exercise}

\begin{exercise}[Posterior mean is optimal]
\label{ch38:ex-bayes}
Let $\Theta \sim \Bet(a,b)$ and, given $\Theta = \theta$, let $D \sim \Bin(n,\theta)$. Find the
posterior $\Theta \given D$ and its mean $\E[\Theta\given D]$. Then argue from the least-squares
characterization \eqref{ch38:l2-problem} that no other function of $D$ achieves smaller expected
squared error against $\Theta$.
\end{exercise}

\begin{exercise}[Why measurability cannot be dropped --- proof]
\label{ch38:ex-measurability}
Show by a counterexample that the averaging property \eqref{ch38:defining-property} \emph{alone},
without the $\mathcal{G}$-measurability requirement, fails to determine a unique object: exhibit two
distinct variables that integrate to $\int_G Y$ over every $G \in \mathcal{G}$ but are not both
$\mathcal{G}$-measurable. Conclude that measurability is doing essential work in pinning down ``what
your information can compute.''
\end{exercise}

\begin{exercise}[Borel--Kolmogorov in coordinates --- proof]
\label{ch38:ex-borel}
On the unit disk with uniform distribution, let $(R,\Theta)$ be polar coordinates. Conditioning on
$\Theta = \theta_0$ (a radius), show the conditional density of $R$ is proportional to $r$ on
$[0,1]$, not uniform. Then re-derive the ``conditional distribution along a radius'' as the limit of
thin angular wedges and confirm the same answer; finally explain in one paragraph how choosing a
different thinning family would change it, mirroring the sphere paradox of Section
\ref{ch38:borel-kolmogorov}.
\end{exercise}
